\documentclass[a4paper]{article}

% Basic packages
% Español (México) con XeLaTeX: fontspec para fuentes Unicode, polyglossia
% para la división silábica y los nombres del documento en la variante mexicana.
\usepackage{fontspec}
\usepackage{polyglossia}
\setdefaultlanguage[variant=mexican]{spanish}
% Los nombres chinos van como «pinyin (original)»: xeCJK da a los caracteres
% Han su propia fuente; sin ella XeLaTeX los omite con un aviso.
% FandolSong no cubre algunos caracteres de nombres (U+73FA); WenQuanYi Zen Hei sí.
\usepackage[AutoFallBack=true]{xeCJK}
\setCJKmainfont{FandolSong-Regular.otf}
\setCJKfallbackfamilyfont{\CJKrmdefault}{WenQuanYi Zen Hei}
% polyglossia no traduce los nombres de listings.
\AtBeginDocument{\providecommand{\lstlistingname}{}\renewcommand{\lstlistingname}{Listado}%
  \providecommand{\lstlistlistingname}{}\renewcommand{\lstlistlistingname}{Índice de listados}}
\usepackage{amsmath, amssymb}
\usepackage{graphicx}
\usepackage[margin=2.5cm]{geometry}
\usepackage[most]{tcolorbox}
\usepackage{etoolbox}
\usepackage{listings}
\usepackage{hyperref}
\usepackage{booktabs}
\usepackage{subcaption}
\usepackage{float}

% Highlight boxes
\newtcolorbox{knowledgebox}[1]{
    enhanced, colback=blue!5!white, colframe=blue!75!black, colbacktitle=blue!75!black,
    coltitle=white, fonttitle=\bfseries, title={#1}, attach boxed title to top left={yshift=-2mm, xshift=2mm},
    boxrule=1pt, sharp corners
}

\newtcolorbox{importantbox}[1]{
    enhanced, colback=yellow!10!white, colframe=yellow!80!black, colbacktitle=yellow!80!black,
    coltitle=black, fonttitle=\bfseries, title={#1}, sharp corners
}

\newtcolorbox{warningbox}[1]{
    enhanced, colback=red!5!white, colframe=red!75!black, colbacktitle=red!75!black,
    coltitle=white, fonttitle=\bfseries, title={#1}, sharp corners
}

% Listings
\lstset{
    language=Python,
    basicstyle=\ttfamily\small,
    keywordstyle=\color{blue},
    stringstyle=\color{red!60!black},
    commentstyle=\color{green!60!black},
    breaklines=true,
    frame=single,
    numbers=left,
    numberstyle=\tiny\color{gray},
    captionpos=b,
    extendedchars=true
}

% Front-page metadata
\newcommand{\notetitle}{CS146S Week 8: Automated UI and App Building}
\newcommand{\noteauthors}{Elaboradas a partir de la clase de Mihail Eric + Gaspar Garcia (Vercel AI Research)}
\newcommand{\notedate}{noviembre de 2025}
\newcommand{\videochannel}{Stanford CS146S}
\newcommand{\videopublishdate}{Fall 2025}
\newcommand{\videoduration}{2 lectures (Mon + Fri)}
\newcommand{\videourl}{https://themodernsoftware.dev}
\newcommand{\videocoverpath}{}
\newcommand{\repourl}{https://github.com/hqhq1025/ai-course-notes}


% Footer with repo info
\usepackage{fancyhdr}
\pagestyle{fancy}
\fancyhf{}
\fancyfoot[C]{\thepage}
\fancyfoot[R]{\footnotesize\color{gray}\href{\repourl}{AI Course Notes}}
\renewcommand{\headrulewidth}{0pt}
\renewcommand{\footrulewidth}{0pt}

\begin{document}

\begin{titlepage}
\centering
{\Large Notas del curso\par}
\vspace{1.2cm}
{\huge\bfseries \notetitle\par}
\vspace{0.8cm}
{\large \noteauthors\par}
\vspace{0.3cm}
{\large \notedate\par}
\vspace{1.2cm}

\ifdefempty{\videocoverpath}{
\vspace{3cm}
}{
\includegraphics[width=0.82\textwidth,height=0.45\textheight,keepaspectratio]{\videocoverpath}\par
}

\vfill
\begin{tcolorbox}[width=0.9\textwidth, colback=black!2!white, colframe=black!60, sharp corners]
\textbf{Curso}: \videochannel\par
\textbf{Semestre}: \videopublishdate\par
\textbf{Clase}: \videoduration\par
\textbf{Sitio web}: \href{\videourl}{\nolinkurl{\videourl}}\par
\textbf{Página del proyecto}: \href{\repourl}{\nolinkurl{\repourl}}\par
\end{tcolorbox}
\end{titlepage}

\tableofcontents
\newpage

%% --- Inicio del contenido --- %%

\section{La evolución del desarrollo de aplicaciones Web}

\subsection{La pila tecnológica Web tradicional}

El desarrollo de SaaS moderno centrado en aplicaciones Web tiene ya más de veinte años. La pila tecnológica tradicional ha pasado por varias generaciones:

\begin{itemize}
    \item \textbf{LAMP}: Linux + Apache + MySQL + PHP
    \item \textbf{MERN}: MongoDB + Express + React + Node.js
    \item \textbf{JAM}: JavaScript + APIs + Markup
    \item \textbf{Serverless}: arquitectura sin servidor
\end{itemize}

Sin importar la pila tecnológica, la arquitectura sigue el patrón clásico de tres capas: cliente (client-side) $\rightarrow$ servidor (server-side) $\rightarrow$ base de datos (database).

\subsection{La primera generación de low-code/no-code}

\begin{itemize}
    \item Herramientas visuales de construcción de sitios como Wix, Squarespace, Webflow
    \item Redujeron la barrera de entrada para construir sitios, pero con capacidad de personalización limitada
\end{itemize}

\subsection{La construcción de aplicaciones en la era de la IA}

\begin{importantbox}{De escribir código a describir necesidades}
El cambio central de las herramientas de construcción de aplicaciones con IA: de «escribir código para construir aplicaciones» a «describir la necesidad en lenguaje natural, y la IA genera la aplicación». Esto es más amigable con el usuario que las plataformas no-code de primera generación, y a la vez conserva una capacidad de personalización más fuerte.
\end{importantbox}

La nueva generación de herramientas de construcción de aplicaciones con IA:
\begin{itemize}
    \item \textbf{Lovable}: generación de aplicaciones full-stack impulsada por IA
    \item \textbf{Replit}: IDE en línea + programación con IA
    \item \textbf{Vercel v0}: lenguaje natural $\rightarrow$ componentes de UI en React
    \item \textbf{Base44}: low-code + IA
    \item \textbf{Cursor / Claude Code}: herramientas de codificación con IA de propósito general que también sirven para construir aplicaciones
\end{itemize}

\begin{knowledgebox}{El cambio en el papel del ingeniero}
Las herramientas de construcción de aplicaciones con IA están volviendo a los ingenieros más \textbf{interfuncionales} (cross-functional): quienes desarrollan ahora también asumen responsabilidades de diseño (design) y de gestión de producto (product management). Al mismo tiempo, las personas sin formación técnica también pueden construir aplicaciones: «Anyone can build an app now.»
\end{knowledgebox}

\subsection{Resumen de la sección}

El desarrollo de aplicaciones Web pasó de LAMP a Serverless, y de ahí a la construcción de aplicaciones con IA; cada generación redujo la barrera de entrada y aumentó la eficiencia. El cambio central de la era de la IA es que el lenguaje natural se convierte en el «lenguaje de programación».

\section{Principios arquitectónicos del App Builder}

\subsection{Arquitectura básica}

La arquitectura central del app builder normalmente se basa en la tecnología WebContainer: ejecuta un entorno de desarrollo completo dentro del navegador, con soporte para edición de código, compilación y vista previa.

\subsection{El motor de v0 en detalle}

Gaspar Garcia, Head of AI Research de Vercel v0, describió en detalle la arquitectura interna de v0:

\subsubsection{Flujo de cinco pasos}

\begin{enumerate}
    \item \textbf{Intent Understanding}: el procesamiento de lenguaje natural extrae los requisitos estructurados: componentes, diseño de la interfaz, interacciones, flujo de datos
    \item \textbf{Context Assembly}: recopila la información relevante: tokens del sistema de diseño, componentes existentes, esquema de la API, requisitos de accesibilidad
    \item \textbf{Code Generation}: el LLM genera componentes de React con calidad de producción (tipos de TypeScript, diseño responsivo, mejores prácticas)
    \item \textbf{Validation \& Testing}: verificaciones automatizadas del cumplimiento de accesibilidad, el comportamiento responsivo y la calidad del código
    \item \textbf{Human-in-the-Loop}: los ingenieros revisan y optimizan, y la retroalimentación se usa para entrenar el modelo
\end{enumerate}

\begin{importantbox}{El Agent es un Workflow}
Gaspar enfatizó una idea central de diseño: un Agent es, en esencia, un workflow (flujo de trabajo). La ``inteligencia'' de v0 no proviene de una sola llamada a un modelo grande, sino de un flujo de trabajo de varios pasos cuidadosamente diseñado. Cada paso tiene entradas, salidas y criterios de validación bien definidos.
\end{importantbox}

\subsection{Limitaciones del modelo de procesamiento}

\subsubsection{Stream Manipulation}

v0 creó un subsistema que puede interceptar y modificar el flujo de salida del LLM antes de que llegue al usuario. Esto resuelve el problema de que ``a veces ajustar solo el prompt no basta'': la corrección se hace directamente a nivel de la salida.

\subsubsection{Entrenamiento y ajuste fino}

Los laboratorios de modelos de frontera no crean canalizaciones de generación de código especialmente diseñadas para aplicaciones web de Next.js 16. v0 construyó su propia composite model family mediante canalizaciones propias de entrenamiento y fine-tuning, para asegurar que el código generado cumpla con las especificaciones más recientes del framework.

\begin{knowledgebox}{Composite Model Family}
v0 no depende de un solo modelo de propósito general, sino que usa un conjunto de modelos ajustados de forma específica y combinados entre sí (composite model family). Distintos modelos se encargan de distintas subtareas, y su combinación logra un resultado conjunto superior al de un solo modelo. Ver el blog de Vercel: \url{https://vercel.com/blog/v0-composite-model-family}
\end{knowledgebox}

\subsection{Resumen de la sección}

El núcleo de los App Builders modernos es un workflow de varios pasos cuidadosamente diseñado, en lugar de una sola llamada al LLM. El stream manipulation y la composite model family son los mecanismos de ingeniería clave para resolver las limitaciones de los modelos de propósito general.
\section{Límites y perspectivas de la construcción de aplicaciones de IA}

\subsection{Límites actuales}

\begin{warningbox}{``Everything is awesome until it breaks''}
El mayor reto de la construcción de aplicaciones de IA: cuando todo funciona bien la experiencia es excelente, pero en cuanto algo falla se vuelve al punto de partida: el desarrollador todavía necesita entender el código subyacente para depurarlo y corregirlo.
\end{warningbox}

\begin{itemize}
    \item \textbf{El prompt es solo una sugerencia}: no todo usuario lo entiende así
    \item \textbf{Seguridad}: históricamente han surgido problemas de seguridad
    \item \textbf{Homogeneización}: ¿cómo evitar que todas las aplicaciones generadas por IA se vean idénticas?
    \item \textbf{Límite de complejidad}: ¿qué tan complejas pueden llegar a ser estas aplicaciones?
\end{itemize}

\subsection{Impacto real}

Las herramientas de construcción de aplicaciones de IA están teniendo impacto en tres niveles:

\begin{itemize}
    \item \textbf{Startup Velocity}: las empresas emergentes lanzan un MVP en días en lugar de meses
    \item \textbf{Enterprise Modernization}: las organizaciones grandes validan primero con prototipos de IA antes de invertir recursos de ingeniería
    \item \textbf{Agency Efficiency}: las agencias de diseño entregan prototipos interactivos horas después de la reunión de arranque
\end{itemize}

\subsection{Criterios del equipo para elegir un App Builder}

Si el contenido del curso se traduce en una decisión de ingeniería, se pueden usar tres preguntas para evaluar si un App Builder de IA en verdad es adecuado para el equipo:

\begin{enumerate}
    \item \textbf{¿Permite iterar de manera continua y no solo generar una vez?} Una herramienta que en verdad aporte valor debe admitir múltiples rondas de modificación, verificación y reversión, en lugar de terminar tras una sola generación.
    \item \textbf{¿Respeta la pila tecnológica existente?} El equipo necesita una herramienta que produzca código que pueda integrarse al React/Next.js/sistema de diseño ya existente, no un prototipo aislado.
    \item \textbf{¿Permite bajar al nivel del código subyacente con rapidez cuando algo falla?} Si la herramienta se vuelve completamente indepurable en cuanto falla, solo sirve como generador de demostraciones, no como herramienta de producción.
\end{enumerate}

\begin{warningbox}{La velocidad del prototipo no es el único indicador}
Los App Builder de IA suelen venderse con el atractivo de «generar una aplicación en minutos», pero lo que en verdad determina el valor a largo plazo es el costo de modificación, el costo de depuración y la compatibilidad con los sistemas de ingeniería ya existentes. Una generación rápida con un costo de mantenimiento posterior alto termina generando una nueva deuda técnica.
\end{warningbox}

\subsection{Resumen de la sección}

Las herramientas de construcción de aplicaciones de IA han reducido de manera notable el tiempo entre la idea y el prototipo, pero el problema de que «en cuanto algo falla se vuelve al punto de partida» nos recuerda que entender la tecnología subyacente sigue siendo indispensable.

\section{De la demo a producción: el verdadero umbral del App Builder}

\subsection{Por qué el éxito del prototipo no equivale al éxito del producto}

Uno de los puntos que más se pasa por alto en el curso es que aquello en lo que el App Builder de IA es más hábil es en hacer visible con rapidez una «idea de interfaz»; pero el trabajo real de conversión en producto sigue existiendo en gran medida.

\begin{itemize}
    \item En la etapa de prototipo basta con hacer que la página y la interacción funcionen
    \item La etapa de producto exige permisos, gestión de estado, manejo de errores, monitoreo, pruebas y mantenibilidad
    \item El lugar donde el equipo realmente invierte tiempo no suele ser generar la primera página de la interfaz, sino conectar el resultado generado con los sistemas de negocio reales
\end{itemize}

\begin{warningbox}{Entre más rápida sea la primera versión generada por la IA, más fácil es subestimar el costo de integración posterior}
Muchos equipos, al ver en cuestión de minutos una «página que funciona», caen en la ilusión de que ya se completó el 80\% del trabajo. Pero en un sistema de producción, lo que en realidad determina el costo suele ser el 20\% restante: el contrato de datos, los límites de error, el modelo de permisos, la revisión y la reversión.
\end{warningbox}

\subsection{El sistema de diseño y el sistema de ingeniería deben volver a conectarse}

Para que el App Builder entre de verdad al flujo de trabajo de un equipo, por lo general debe conectar dos sistemas:

\begin{enumerate}
    \item \textbf{El sistema de diseño}: tokens de color, especificaciones de componentes, tipografía, restricciones de interacción
    \item \textbf{El sistema de ingeniería}: versión del framework, estructura de directorios, patrones de obtención de datos, pruebas y despliegue
\end{enumerate}

\begin{table}[H]
\centering
\begin{tabular}{p{0.2\textwidth} p{0.28\textwidth} p{0.28\textwidth}}
\toprule
Capa del sistema & Qué pasa si falta & Efecto sobre el App Builder \\
\midrule
Sistema de diseño & La página funciona pero la apariencia visual no es consistente & Cada generación parece un producto distinto ensamblado \\
Especificación de componentes & El código es difícil de reutilizar y la interacción es inconsistente & La IA no puede reutilizar de forma estable las piezas ya existentes \\
Contrato de datos & La interfaz queda desconectada del backend real & Solo se puede permanecer en la capa de la demo simulada \\
Pruebas y despliegue & El resultado generado no puede publicarse de forma estable & El equipo no se atreve a conectar el resultado de la IA al flujo principal \\
\bottomrule
\end{tabular}
\caption{Capas del sistema que deben completarse para pasar de herramienta de prototipo a herramienta de producción}
\end{table}

\begin{knowledgebox}{El mejor App Builder no sustituye al sistema de ingeniería, sino que acorta la distancia para conectarse a él}
Las herramientas verdaderamente destacadas no fingen que la complejidad de la ingeniería no existe, sino que procuran alinear el resultado generado con el sistema de diseño, el sistema de tipos, el proceso de pruebas y la forma de despliegue que el equipo ya tiene. Así, la IA es un acelerador y no una nueva isla aislada.
\end{knowledgebox}

\subsection{Resumen de la sección}

El valor central del App Builder está en convertir con rapidez una idea en una interfaz que se puede discutir, pero para cambiar de verdad el flujo de trabajo de un equipo debe atravesar umbrales más duros: el sistema de diseño, el contrato de datos, las pruebas y el despliegue.

\section{Síntesis y ampliación}

La semana 8 mostró cómo la IA transforma la construcción de aplicaciones: de la pila LAMP tradicional al paradigma impulsado por IA de lenguaje natural $\rightarrow$ aplicación. El análisis detallado de la arquitectura de v0 reveló que detrás de la «magia» del App Builder hay un workflow y un composite model cuidadosamente diseñados.

\subsection{Puntos clave}

\begin{itemize}
    \item El desarrollo web evolucionó de LAMP $\rightarrow$ MERN $\rightarrow$ JAM $\rightarrow$ Serverless $\rightarrow$ AI App Builder
    \item La IA hace que los ingenieros sean más multifuncionales, y las personas sin conocimientos técnicos también pueden construir aplicaciones
    \item Arquitectura de v0: Intent $\rightarrow$ Context $\rightarrow$ Code Gen $\rightarrow$ Validation $\rightarrow$ Human Review
    \item «Agent es Workflow» — lo esencial no es una única llamada al modelo, sino una cadena de procesamiento de varios pasos
    \item Stream manipulation y composite model family resuelven las limitaciones de los modelos de propósito general
    \item La mayor limitación: cuando ocurre un error, todavía se necesita que una persona entienda el código subyacente para depurarlo
\end{itemize}

\subsection{Hoja de ruta de implementación}

\begin{table}[H]
\centering
\begin{tabular}{p{0.18\textwidth} p{0.26\textwidth} p{0.3\textwidth}}
\toprule
Etapa & Objetivo principal & Práctica típica \\
\midrule
Validación de prototipo & Validar rápidamente si la idea es viable & Generar con lenguaje natural páginas, formularios y el esqueleto de las interacciones \\
Integración con el sistema & Conectar datos reales y el sistema de diseño & Conectar APIs, completar tipos y alinear las especificaciones de los componentes \\
Puesta en producción & Garantizar un mantenimiento estable y la colaboración de varias personas & Agregar pruebas, hacer review, y consolidar el workflow y los mecanismos de reversión \\
\bottomrule
\end{tabular}
\caption{Ruta típica de adopción de AI App Builder}
\end{table}

\subsection{Lista de acciones para el equipo}

\begin{enumerate}
    \item Elegir primero una herramienta interna de bajo riesgo o un proyecto de prototipo para hacer un piloto, en lugar de aplicarlo directamente a las páginas centrales del negocio
    \item Definir de antemano el sistema de diseño y los límites de los componentes, para que el resultado generado se ajuste de manera estable a las normas de ingeniería existentes
    \item Incorporar el resultado generado por la IA al proceso normal de code review, pruebas y reversión, en lugar de tratarlo como un proceso excepcional
\end{enumerate}

\subsection{Lecturas adicionales}

\begin{itemize}
    \item Vercel v0: \url{https://v0.dev} (gratis para estudiantes: \url{https://v0.app/students})
    \item Vercel AI SDK: \url{https://ai-sdk.dev/docs/agents/overview}
    \item Publicación de blog sobre Composite Model Family: \url{https://vercel.com/blog/v0-composite-model-family}
    \item Lovable: \url{https://lovable.dev}
    \item Bolt.new (StackBlitz): \url{https://bolt.new}
    \item Replit: \url{https://replit.com}
\end{itemize}

%% --- Fin del contenido --- %%

\end{document}
